\worksheettitle
  {Цифры, числа и разряды}
  {\modulemark{M01} Версия учителя}

\begin{teacherbox}[Паспорт модуля]
  \textbf{Результат:} ученик определяет разряд цифры и называет её значение.

  \textbf{Предварительно:} различает цифры в записи числа и знает названия
  единиц, десятков и сотен.

  \textbf{Ориентир:} 15--20 минут. Не переходите дальше только потому, что
  время истекло: используйте контрольную точку.
\end{teacherbox}

\begin{checkpointbox}[Вход: ответы и решение]
  1) Три цифры: \(5,0,5\). \quad 2) Цифра \(8\).

  Если ученик не считает ноль цифрой, проговорите запись \(505\) по позициям.
  Если считает разряды слева, быстро разберите трёхзначное число справа
  налево. После одной устной опоры можно начинать M01.
\end{checkpointbox}

\section{Открытие способа}

Заполненная таблица:

\begin{center}
\renewcommand{\arraystretch}{1.38}
\begin{tabularx}{0.9\linewidth}{@{}>{\centering\arraybackslash}p{0.18\linewidth}>{\centering\arraybackslash}X>{\centering\arraybackslash}X@{}}
  \toprule
  \rowcolor{TableHeader}
  Число & Разряд цифры \(6\) & Значение цифры \(6\) \\
  \midrule
  \(6\) & единицы & \(6\) \\
  \(60\) & десятки & \(60\) \\
  \(600\) & сотни & \(600\) \\
  \(6\,000\) & единицы тысяч & \(6\,000\) \\
  \bottomrule
\end{tabularx}
\end{center}

\begin{teacherbox}[Вопросы учителя]
  \begin{itemize}
    \item Какая цифра повторяется во всех записях?
    \item Что изменяется, хотя сама цифра остаётся той же?
    \item Чем отличается ответ «цифра 6» от ответа «значение 600»?
  \end{itemize}
  Главный вывод: значение цифры зависит от её места, то есть от разряда.
\end{teacherbox}

\section{Практика: ответы}

Для числа \(72\,649\):

\begin{center}
\renewcommand{\arraystretch}{1.32}
\begin{tabular}{c|ccccc}
  \toprule
  цифра & 7 & 2 & 6 & 4 & 9 \\
  \midrule
  разряд & дес. тыс. & ед. тыс. & сотни & десятки & единицы \\
  значение & 70 000 & 2 000 & 600 & 40 & 9 \\
  \bottomrule
\end{tabular}
\end{center}

Самостоятельная проба:
\begin{enumerate}
  \item сотни, \(300\);
  \item единицы тысяч, \(6\,000\);
  \item \(50\,000\) и \(50\).
\end{enumerate}

\begin{teacherbox}[Наблюдайте не только ответ]
  Просите ученика говорить двумя фразами: «Цифра стоит в разряде…» и
  «Поэтому её значение…». Ответ только \(6\) на вопрос о значении показывает,
  что различие ещё не сформировано.
\end{teacherbox}

\section{Контрольная точка и маршрут}

Для числа \(44\,307\): первая цифра \(4\) имеет значение \(40\,000\),
вторая --- \(4\,000\). Они занимают разные разряды.

\begin{resultbox}[Решение о переходе]
  \begin{itemize}
    \item \textbf{Переход к M02:} оба значения названы правильно, ученик
      объясняет зависимость от разряда.
    \item \textbf{M01-R:} хотя бы одно значение заменено самой цифрой,
      разряды считаются слева или объяснение отсутствует.
    \item \textbf{M01\(\star\):} базовый результат достигнут без опоры; можно дать
      задание на три одинаковые цифры в одном числе.
  \end{itemize}
\end{resultbox}

\begin{teacherbox}[После коррекции]
  Не повторяйте число \(44\,307\). Дайте новую проверку: «Назовите значения
  двух цифр \(8\) в числе \(81\,804\)». Ожидается \(80\,000\) и \(800\).
\end{teacherbox}
